%=====================================================================
\chapter{Data Engineering at Scale}\label{ch:dataeng}
%=====================================================================
\locator{5}{Part V --- Data Science in Systems Context}

\begin{outcomes}
By the end of this chapter you will be able to:
\begin{tight}
  \item \textbf{Describe} a data pipeline and distinguish ETL from ELT.
  \item \textbf{Compare} data warehouses, lakes and lakehouses, and choose one for a stated need.
  \item \textbf{Explain} the five V's of big data and what each demands technically.
  \item \textbf{Describe} how partitioning and parallelism make large computations tractable.
  \item \textbf{Specify} data quality checks that run automatically in a pipeline.
\end{tight}
\end{outcomes}

\begin{prereq}
\chref{data} (quality dimensions), \chref{wrangling} (the transformations a
pipeline automates) and \chref{timeseries} (batch versus stream).
\end{prereq}

\begin{keyterms}
data pipeline $\bullet$ ETL / ELT $\bullet$ orchestration $\bullet$ DAG $\bullet$
idempotence $\bullet$ data warehouse $\bullet$ data lake $\bullet$ lakehouse
$\bullet$ schema-on-read / schema-on-write $\bullet$ the five V's $\bullet$
partitioning $\bullet$ MapReduce $\bullet$ Spark $\bullet$ data contract
$\bullet$ lineage
\end{keyterms}

\section{Why Data Engineering Exists}

\begin{conceptbox}[The Uncomfortable Ratio]
A model that runs once in a notebook needs no engineering. A model that must run
every morning on yesterday's data, for the next three years, needs a great deal of
it. Most data science that reaches production is 20\% modelling and 80\% making
data arrive reliably --- and the 80\% is what this chapter is about.
\end{conceptbox}

\section{Pipelines}

\begin{definitionbox}[Data Pipeline]
An automated, scheduled sequence of steps that moves data from source systems to a
form analysts and models can use. Steps are arranged as a \term{directed acyclic
graph} (DAG): each task declares what it depends on, and the orchestrator works out
the order and what can run in parallel.
\end{definitionbox}

\dsfig{%
\begin{tikzpicture}[font=\scriptsize,
  src/.style={rounded corners=3pt, draw=secondary, fill=secondary!9, text=secondary!55!black,
              text width=1.9cm, align=center, font=\tiny\bfseries, minimum height=0.72cm},
  t/.style={rounded corners=3pt, draw=#1, fill=#1!9, text=#1!50!black,
            text width=2.0cm, align=center, font=\tiny\bfseries, minimum height=0.72cm},
  arr/.style={-{Stealth[length=1.8mm]}, gray!60, line width=0.8pt}]

\node[src] (s1) at (0,1.5)   {registry DB};
\node[src] (s2) at (0,0.55)  {VLE API};
\node[src] (s3) at (0,-0.4)  {swipe logs};
\node[src] (s4) at (0,-1.35) {sensor stream};

\node[t=primary]  (ing) at (2.9,0.07)  {ingest\\{\tiny raw, unchanged}};
\node[t=goldhl]   (val) at (5.6,0.07)  {validate\\{\tiny contracts}};
\node[t=tealhl]   (tr1) at (8.3,0.85)  {clean \&\\conform};
\node[t=tealhl]   (tr2) at (8.3,-0.7)  {aggregate\\{\tiny to grain}};
\node[t=greenhl]  (feat) at (11.0,0.07) {feature\\store};
\node[t=purplehl] (out) at (13.4,0.07)  {model /\\dashboard};

\foreach \s in {s1,s2,s3,s4}{\draw[arr] (\s) -- (ing);}
\draw[arr] (ing) -- (val);
\draw[arr] (val) -- (tr1);
\draw[arr] (val) -- (tr2);
\draw[arr] (tr1) -- (feat);
\draw[arr] (tr2) -- (feat);
\draw[arr] (feat) -- (out);

% quarantine branch
\node[t=accent] (q) at (5.6,-1.75) {quarantine};
\draw[arr, accent] (val) -- node[right, font=\tiny, text=accent!75!black, pos=0.45] {fails} (q);

\node[anchor=west, align=left, text width=13.6cm, font=\tiny, text=gray!55!black]
     at (-1.2,-2.9)
     {\textbf{Three properties that separate a pipeline from a script:}
      it is \textbf{idempotent} (re-running Tuesday's job produces Tuesday's result, not a
      duplicate); it \textbf{fails loudly} (bad data is quarantined and alerted, never
      silently averaged in); and it records \textbf{lineage} (for any number on the
      dashboard you can say which source rows produced it).};
\end{tikzpicture}}%
{A data pipeline as a DAG. The quarantine branch is the part beginners omit and the part
that saves projects: data that fails validation must stop, not flow onwards.}{pipeline-dag}

\begin{definitionbox}[ETL versus ELT]
\term{ETL}: Extract, Transform, then Load --- transform before storing, so the
warehouse holds only clean, modelled data. \term{ELT}: Extract, Load, then
Transform --- land the raw data first and transform inside the destination system.
ELT dominates modern practice because storage is cheap and keeping the raw data
lets you re-derive everything when requirements change.
\end{definitionbox}

\section{Where Data Lives}

\dsfig{%
\begin{tikzpicture}[font=\scriptsize,
  hd/.style={rounded corners=3pt, fill=#1, text=white, font=\tiny\bfseries,
             text width=3.9cm, align=center, minimum height=0.55cm},
  b/.style={rounded corners=3pt, draw=#1, fill=#1!7, text width=3.9cm, align=left,
            font=\tiny, inner sep=5pt, minimum height=3.5cm}]

\node[hd=secondary] at (0,0)    {DATA WAREHOUSE};
\node[b=secondary]  at (0,-2.15){%
  \textbf{Schema-on-write:} structure is fixed before loading.\\[4pt]
  \textbf{Holds:} cleaned, modelled tables\\
  \textbf{Query:} SQL, fast\\
  \textbf{Users:} analysts, BI tools\\[4pt]
  {\color{greenhl!45!black}\textbf{+} fast, governed, trustworthy}\\
  {\color{accent}\textbf{--} rigid; adding a source is a project}\\[4pt]
  \emph{Snowflake, BigQuery, Redshift}};

\node[hd=goldhl] at (4.5,0)    {DATA LAKE};
\node[b=goldhl]  at (4.5,-2.15){%
  \textbf{Schema-on-read:} store anything, interpret later.\\[4pt]
  \textbf{Holds:} raw files of any type\\
  \textbf{Query:} engines over files\\
  \textbf{Users:} data scientists, engineers\\[4pt]
  {\color{greenhl!45!black}\textbf{+} cheap, flexible, keeps everything}\\
  {\color{accent}\textbf{--} becomes a ``data swamp'' without catalogue and governance}\\[4pt]
  \emph{S3, ADLS, GCS}};

\node[hd=greenhl] at (9.0,0)    {LAKEHOUSE};
\node[b=greenhl]  at (9.0,-2.15){%
  Lake storage plus warehouse guarantees.\\[4pt]
  \textbf{Holds:} raw \emph{and} curated\\
  \textbf{Query:} SQL over open formats\\
  \textbf{Users:} both audiences\\[4pt]
  {\color{greenhl!45!black}\textbf{+} transactions, schema evolution, one copy of the data}\\
  {\color{accent}\textbf{--} newer; more moving parts}\\[4pt]
  \emph{Delta Lake, Iceberg}};

\node[anchor=north, align=left, text width=13.2cm, font=\scriptsize, text=gray!55!black]
     at (4.5,-4.2)
     {\textbf{How to choose.} Known questions, structured sources and BI reporting
      $\Rightarrow$ warehouse. Unknown future questions, unstructured or sensor data
      $\Rightarrow$ lake. Both, and you can afford the extra moving parts
      $\Rightarrow$ lakehouse.};
\end{tikzpicture}}%
{Warehouse, lake and lakehouse. The real distinction is \emph{when} you commit to a
schema --- before loading, or at query time --- and that choice determines how expensive
it is to change your mind later.}{warehouse-lake}

\section{Big Data: The Five V's}

\rowcolors{2}{white}{lightbg}
\begin{longtable}{@{}p{2.3cm}p{4.3cm}p{7.4cm}@{}}
\toprule
\rowcolor{primary!15}
\textbf{V} & \textbf{Meaning} & \textbf{What it forces you to do} \\
\midrule
\endhead
Volume & More data than one machine holds & Partition across machines; move the computation to the data, not the reverse \\
Velocity & Arriving faster than batch can absorb & Stream processing and windowing (\chref{timeseries}) \\
Variety & Structured, semi-structured and unstructured together & Schema-on-read; a lake rather than a warehouse \\
Veracity & Quality varies and cannot be assumed & Automated validation, quarantine, lineage \\
Value & Most of it is never used & Prioritise ruthlessly; storage is cheap but attention is not \\
\bottomrule
\end{longtable}

\begin{alertbox}[``Big Data'' Is Rarer Than Claimed]
A dataset is big when it does not fit on one machine. A year of records for
2{,}000 students is a few hundred megabytes and belongs in pandas, not in Spark.
Reaching for distributed tooling on data that fits in memory adds days of
complexity for negative benefit. Measure first.
\end{alertbox}

\section{Parallelism and Partitioning}

\begin{definitionbox}[Partitioning]
Split data into independent chunks --- by date, by device, by region --- so that
each can be processed on a different machine simultaneously, and so that a query
for one day reads one partition instead of the whole dataset.
\end{definitionbox}

\dsfig{%
\begin{tikzpicture}[font=\scriptsize,
  blk/.style={rounded corners=2pt, draw=#1, fill=#1!12, text=#1!50!black,
              minimum width=1.55cm, minimum height=0.55cm, font=\tiny, align=center},
  arr/.style={-{Stealth[length=1.8mm]}, gray!60, line width=0.8pt}]

\node[blk=primary, minimum width=2.2cm] (in) at (0,0) {400 GB of\\telemetry};

\foreach \i/\lbl in {1/{node 1\\Jan}, 2/{node 2\\Feb}, 3/{node 3\\Mar}, 4/{node 4\\Apr}}{
  \pgfmathsetmacro{\yy}{1.65-(\i-1)*1.1}
  \node[blk=secondary, minimum height=0.8cm] (m\i) at (3.3,\yy) {\lbl};
  \draw[arr] (in) -- (m\i);
  \node[blk=goldhl, minimum height=0.8cm] (r\i) at (6.2,\yy) {partial\\result};
  \draw[arr] (m\i) -- (r\i);}

\node[blk=greenhl, minimum width=2.2cm, minimum height=1.0cm] (out) at (9.2,0.0)
     {combined\\result};
\foreach \i in {1,2,3,4}{\draw[arr] (r\i) -- (out);}

\node[font=\tiny\bfseries, text=secondary!55!black] at (3.3,2.6) {MAP};
\node[font=\tiny\bfseries, text=goldhl!50!black] at (6.2,2.6) {(shuffle)};
\node[font=\tiny\bfseries, text=greenhl!45!black] at (9.2,2.6) {REDUCE};

\node[anchor=west, align=left, text width=4.2cm, font=\tiny, text=gray!55!black] at (10.8,0.0)
     {\textbf{Why this works:} the four months are independent, so the map step needs no
      communication. Operations that \emph{do} need communication --- a join, a global
      sort --- force a shuffle across the network, which is the expensive part. Design
      partitions so the common query touches one.};

\node[anchor=west, align=left, text width=13.6cm, font=\tiny, text=accent!75!black]
     at (-1.2,-2.55)
     {\textbf{Partition by what you filter on.} If every query asks for one device's last
      week, partition by date and device. Partitioning by something you never filter on
      buys nothing and multiplies your file count --- the ``small files problem'' that
      cripples lakes.};
\end{tikzpicture}}%
{Map--reduce over partitioned data. Spark is a modern, in-memory implementation of this
idea; the principle --- split, compute independently, combine --- is what makes data larger
than one machine tractable.}{mapreduce}

\section{Data Contracts and Automated Quality}

\begin{definitionbox}[Data Contract]
An explicit, versioned agreement between a data producer and its consumers: which
fields exist, their types, their permitted ranges, their nullability, and the
freshness guarantee. Enforced automatically at ingestion, it converts a silent
breakage into a loud, attributable failure.
\end{definitionbox}

\begin{worked}[Writing Checks That Would Have Caught Real Failures]
Each check below corresponds to a failure from the audit in \chref{data}.

\smallskip
\begin{center}
\begin{tabular}{@{}p{3.5cm}p{5.2cm}p{5.0cm}@{}}
\toprule
\textbf{Dimension} & \textbf{Automated check} & \textbf{Action on failure} \\
\midrule
Completeness & \texttt{null\_rate(prior\_qual) < 0.40} & warn; block if $>0.6$ \\
Validity & \texttt{final\_grade IN (A,B,C,D,F)} & quarantine the rows \\
Accuracy & \texttt{age BETWEEN 15 AND 90} & quarantine, alert owner \\
Uniqueness & \texttt{count(*) = count(DISTINCT student\_id, term)} & block the run \\
Timeliness & \texttt{max(loaded\_at) > now() - 26h} & page on-call \\
Consistency & \texttt{|registry\_count $-$ vle\_count| < 50} & warn \\
Volume & \texttt{row\_count} within $\pm30\%$ of the 7-day median & block the run \\
\bottomrule
\end{tabular}
\end{center}

\smallskip
\textbf{The last check is the most valuable and the least obvious.} A source that
silently starts returning 10\% of its usual rows breaks no type and violates no
range --- every individual row is valid. Only the \emph{row count} reveals it, and
without that check the model quietly trains on a biased tenth of the population.
Volume anomaly detection is the cheapest high-value check in any pipeline.
\end{worked}

%---------------------------------------------------------------------
\begin{fourlenses}{Data Engineering}
\lensLA{Sources are the registry (nightly batch), the VLE (API, hourly) and swipe
readers (near real time). Grain conflicts are the recurring problem: registry is
one row per enrolment, VLE is one row per click. Land both raw, then build a
conformed weekly student table. \textbf{Governance is unusually strict:} this is
personal data about identifiable young people, so lineage is not merely good
practice but a requirement (\chref{ethics}).}

\lensPM{Volumes are trivial --- everything fits in a spreadsheet --- so the
engineering problem is not scale but \emph{integration}. Ticket tracker, git,
timesheets and CI logs each identify the same work item differently, and reconciling
those identifiers is most of the work. A data contract with each tool's owner is
worth more here than any distributed technology.}

\lensIOT{The genuine big-data case in this course: 400 devices $\times$ 6 sensors
$\times$ 1 Hz is 200 million rows a day. Partition by date and device, store in a
columnar format, and pre-aggregate to minute and hour grains --- almost no query
needs raw 1 Hz data. Expect duplicates from at-least-once delivery and out-of-order
arrival from network delay, so ingestion must be idempotent and windowing must use
watermarks (\chref{timeseries}).}

\lensSEC{Log volumes are enormous and retention is often mandated by regulation, so
the lake pattern dominates: land everything raw, index selectively. Two engineering
requirements are specific to this domain --- logs must be \textbf{tamper-evident},
because an attacker's first act is to alter them, and lineage must be strong enough
to stand up in an investigation. Ingestion latency is itself a security control: a
detection pipeline with a six-hour lag cannot stop anything.}
\end{fourlenses}

%---------------------------------------------------------------------
\begin{lab}[A Pipeline Task with Contracts]
\begin{lstlisting}[language=Python]
import pandas as pd
from datetime import datetime, timedelta

class DataContractViolation(Exception):
    pass

CONTRACT = {
    "required":  ["student_id", "term", "attendance_pct", "loaded_at"],
    "ranges":    {"attendance_pct": (0, 100), "age": (15, 90)},
    "allowed":   {"final_grade": {"A", "B", "C", "D", "F"}},
    "unique_on": ["student_id", "term"],
    "max_null":  {"prior_qualification": 0.40},
    "freshness_hours": 26,
}

def validate(df, contract=CONTRACT):
    """Returns (clean_rows, quarantined_rows). Blocks on structural failures."""
    missing = set(contract["required"]) - set(df.columns)
    if missing:
        raise DataContractViolation(f"missing required columns: {missing}")

    if df.duplicated(contract["unique_on"]).any():
        raise DataContractViolation("uniqueness violated -- run blocked")

    age = (datetime.now() - df["loaded_at"].max()).total_seconds() / 3600
    if age > contract["freshness_hours"]:
        raise DataContractViolation(f"stale data: {age:.1f}h old")

    bad = pd.Series(False, index=df.index)
    for col, (lo, hi) in contract["ranges"].items():
        if col in df:
            bad |= ~df[col].between(lo, hi) & df[col].notna()
    for col, allowed in contract["allowed"].items():
        if col in df:
            bad |= ~df[col].isin(allowed) & df[col].notna()

    for col, limit in contract["max_null"].items():
        if col in df and df[col].isna().mean() > limit:
            print(f"WARN: {col} is {df[col].isna().mean():.0%} null (limit {limit:.0%})")

    return df[~bad], df[bad]          # quarantine, never silently drop

def check_volume(df, history_median, tolerance=0.30):
    """The cheapest high-value check: did the source quietly shrink?"""
    ratio = len(df) / history_median
    if abs(ratio - 1) > tolerance:
        raise DataContractViolation(
            f"row count {len(df)} is {ratio:.0%} of the 7-day median -- run blocked")

# clean, quarantined = validate(raw_df)
# check_volume(clean, history_median=1987)
# quarantined.to_parquet(f"quarantine/{datetime.now():%Y-%m-%d}.parquet")
\end{lstlisting}
\end{lab}

%---------------------------------------------------------------------
\begin{checkpoint}
\begin{tightnum}
  \item What does it mean for a pipeline task to be idempotent, and why does it
        matter when a job is re-run after a failure?
  \item Your dataset is 300 MB. Should you use Spark? Justify.
  \item A source silently starts returning 10\% of its usual rows. Which of the six
        quality dimensions catches this, and which check specifically?
\end{tightnum}
\end{checkpoint}

%---------------------------------------------------------------------
\begin{chaptersummary}
\begin{tight}
  \item Pipelines are DAGs of scheduled tasks; they must be idempotent, fail loudly
        and record lineage.
  \item ELT has largely displaced ETL because cheap storage makes keeping the raw
        data worthwhile.
  \item Warehouses fix the schema on write, lakes on read, lakehouses attempt both.
        The choice is about when you commit.
  \item The five V's each impose a technical requirement; veracity is the one most
        often ignored.
  \item Partition by what you filter on. Shuffles across the network are the
        expensive operation.
  \item Data contracts turn silent breakage into an attributable failure, and a
        row-count check is the cheapest valuable test you can write.
\end{tight}
\end{chaptersummary}

\begin{reviewq}
  \item \textbf{(MCQ)} Schema-on-read is characteristic of: (a)~a data warehouse
        (b)~a data lake (c)~ETL (d)~a relational database.
  \item \textbf{(MCQ)} Which operation forces an expensive network shuffle?
        (a)~filtering one partition (b)~a per-row transformation (c)~a join across
        partitions (d)~reading a column.
  \item \textbf{(Short)} Define idempotence and give an example of a non-idempotent
        pipeline step and its consequence.
  \item \textbf{(Short)} Explain why a quarantine branch is preferable to dropping
        invalid rows.
  \item \textbf{(Short)} Why is ingestion latency a security control in a detection
        pipeline?
  \item \textbf{(Applied)} Design the ingestion layer for 400 IoT devices reporting
        6 channels at 1 Hz. Specify the storage format, partitioning scheme,
        aggregation grains, and how you handle duplicates and late arrivals.
  \item \textbf{(Applied)} Write six data-contract checks for a nightly student
        enrolment feed, stating for each whether failure should warn, quarantine or
        block the run --- and justify the three that block.
\end{reviewq}
